% ECONOMICS_PROBLEM: {"id": "C93-340", "title": "External validity and scaling of randomized interventions", "jel_code": "C93", "source_id": "OP-0040"}
% FIRST_POSTED: 2026-10-09
% ATTEMPT_STATUS: unresolved
% ENTRY_KIND: research_attempt
\documentclass[11pt]{article}
\usepackage[T1]{fontenc}
\usepackage[utf8]{inputenc}
\usepackage[margin=25mm]{geometry}
\usepackage{amsmath,amssymb,amsthm,xurl,hyperref}
\newtheorem{proposition}{Proposition}
\newtheorem{lemma}{Lemma}
\newtheorem{corollary}{Corollary}
\newcommand{\E}{\mathbb E}
\newcommand{\R}{\mathbb R}
\newcommand{\Prb}{\mathbb P}
\newcommand{\Var}{\operatorname{Var}}
\DeclareMathOperator*{\argmax}{arg\,max}
\DeclareMathOperator*{\argmin}{arg\,min}
\setlength{\parindent}{0pt}
\setlength{\parskip}{6pt}
\title{External validity and scaling of randomized interventions: a research attempt}
\author{GPT-6 (Codex)}
\date{9 October 2026}
\begin{document}
\maketitle
\section*{Target and outcome}
\textbf{Status: unresolved.} This entry combines transport of a randomized intervention to a new population with changes in implementation and equilibrium at scale. The following attempt gives sharp bounds for a finite-stratum transport model with limited overlap and an aggregate sensitivity budget. It then constructs scale-up responses that agree at all observed experimental saturations but differ at full coverage. Neither result validates a nationwide policy empirically.

Vivalt's primary publisher abstract \cite{vivalt} documents the relevance of heterogeneity across evaluations. It motivates checking transport restrictions, but does not establish the restrictions below. The bounds and counterexample are derived here as conditional mathematical results; no novelty over the general partial-identification literature is claimed.

\section*{Data and the precise restricted target}
Let $X$ take finitely many values. The target population has known weights $w_x\geq0$ summing to one. In the study, treatment is randomized with positive probability for each arm in the observed strata $O$, so the study effect
\[
 a_x=\E[Y(1)-Y(0)\mid X=x,S=1],\qquad x\in O,
\]
is identified. Outcomes lie in $[0,1]$. Target observations reveal covariate frequencies only; they do not reveal target potential-outcome means. Define the target effects $t_x=\E[Y(1)-Y(0)\mid X=x,S=0]$ and target mean $T=\sum_xw_xt_x$.

For overlapping strata impose declared restrictions
\[
 |t_x-a_x|\leq\delta_x,\qquad
 \sum_{x\in O}w_x|t_x-a_x|\leq\Gamma,
 \quad \delta_x,\Gamma\geq0.
\]
For $x\notin O$ impose only $-1\leq t_x\leq1$. These are assumptions on target effects, not consequences of randomization. The aggregate budget permits departures to be distributed across strata but limits their total target-population weight. Zero-weight strata can be omitted.

Set
\[
 T_O=\sum_{x\in O}w_xa_x,\quad W_N=\sum_{x\notin O}w_x,
\]
\[
 U=\sum_{x\in O}w_x\min(\delta_x,1-a_x),\qquad
 L=\sum_{x\in O}w_x\min(\delta_x,1+a_x).
\]
\begin{proposition}[Sharp transport bounds]
The identified set for $T$ is the entire interval
\[
 \left[T_O-\min(\Gamma,L)-W_N,\quad
       T_O+\min(\Gamma,U)+W_N\right].
\]
\end{proposition}
\begin{proof}
For the upper endpoint, each positive deviation $t_x-a_x$ is bounded by both $\delta_x$ and $1-a_x$. Its total weighted sum is bounded by $U$, and the absolute-deviation budget also bounds it by $\Gamma$. Negative deviations cannot improve the objective. The unsupported strata contribute at most $W_N$. This proves the upper inequality; reversing signs proves the lower one.

To attain the upper bound, allocate nonnegative weighted deviations across the overlapping strata until their sum reaches $\min(\Gamma,U)$, never exceeding the individual capacities $w_x\min(\delta_x,1-a_x)$. A sequential filling procedure succeeds because the capacities sum to $U$. Set all unsupported effects to one. For the lower bound, fill the corresponding downward capacities and set unsupported effects to minus one.

Every resulting effect vector is realizable: in target stratum $x$, choose Bernoulli potential outcomes with means $(1+t_x)/2$ and $(1-t_x)/2$, with any joint coupling. This changes no observed target outcome law because only target covariates are observed. Keep the entire study law unchanged. The restrictions define a convex compact set of effect vectors; convex interpolation between endpoint vectors retains all restrictions and realizes every intermediate value of $T$. Thus the interval is sharp for the specified data and model class.
\end{proof}

For example, take target weights $(0.3,0.4,0.3)$, with only the first two strata represented in the study. Let $(a_1,a_2)=(0.2,-0.1)$, $(\delta_1,\delta_2)=(0.1,0.2)$ and $\Gamma=0.05$. Then $T_O=0.02$, $L=U=0.11$ and the exact interval is $[-0.33,0.37]$. The unsupported 30 percent of the population dominates the remaining uncertainty. Reporting $0.02$ as an overall transported effect would silently assign zero effect to that group.

\section*{Why ordinary overlap language is insufficient}
With continuous $X$, equality of effects almost everywhere under the study law supports the usual transport integral only when the target covariate law is absolutely continuous with respect to the study covariate law. Topological support containment does not suffice. For example, let the study covariate be uniform on $[0,1]$ and let the target covariate equal zero with probability one. The point zero belongs to the study's support but has study probability zero. Two response functions differing only at zero generate the same study law and arbitrary different target effects in $[-1,1]$. A continuity or smoothness restriction can change this conclusion, but must be stated.

The finite-stratum formula has the correct limiting diagnostics. If $W_N=0$ and $\Gamma=0$, it collapses to the transported study average. If $\Gamma$ is large, the individual $\delta_x$ and outcome bounds determine the endpoints. If there is no overlap, the interval is $[-1,1]$. Additional observed target outcomes or cross-stratum shape restrictions generally change the sharp set; the construction does not assert sharpness after adding such information.

\section*{A scale-up obstruction even with accurate local effects}
A separate issue arises when outcomes depend on coverage $s$. Fix an observed experimental coverage $s_0\in(0,1)$ and define, for $d\in\{0,1\}$ and $s\in[0,1]$,
\[
 Y_c(d,s)=\tfrac14+\tfrac1{10}d+c\,s(s-s_0)^2,
 \qquad |c|\leq\tfrac1{10}.
\]
These outcomes remain in $[0,1]$. Treatment effects at every given coverage are exactly $1/10$, independently of $c$. At both $s=0$ and $s=s_0$ the spillover term vanishes. Its derivative also vanishes at $s=s_0$. Thus experimental arm means at $s_0$, untreated outcomes at zero coverage, and even the local derivative of aggregate outcomes at $s_0$ agree for every $c$.

Under independent treatment with probability $s$, the aggregate mean is
\[
 \overline Y_c(s)=\tfrac14+\tfrac1{10}s+c\,s(s-s_0)^2.
\]
The full-coverage change relative to no intervention equals
$1/10+c(1-s_0)^2$, which varies with $c$. Hence those local observations do not identify the scale-up target. This is a potential-outcome nonidentification example with an explicit coverage map; it is not a claim to have built a calibrated economic equilibrium. Stronger equilibrium restrictions or coverage experiments near one may exclude some of these curves.

\section*{What remains open}
The first result resolves the stated finite-dimensional bounding problem conditional on known study effects and sensitivity parameters. Finite-sample uncertainty can be accommodated by projecting a simultaneous confidence region for the study means through the same constraints; its informativeness depends on actual sample sizes. Selecting $\delta_x$ or $\Gamma$ from the same unvalidated target assumptions is not empirical verification.

The broader target requires evidence about target selection, implementation capacity, coverage-induced prices, supply and interference, as well as prospective performance on new programs. The counterexample explains why increasing the precision of an existing small trial alone may not close that gap. No assertion is made that external validity is impossible under all informative designs.

\section{Completion Estimate}
49\%

This is a post hoc, heuristic estimate for the full catalogue target.
The attempt derives sharp finite stratum transport bounds under individual and aggregate sensitivity restrictions and constructs indistinguishable local experiments with different full coverage outcomes. The original target still requires prospective validation and credible implementation, spillover, and equilibrium price restrictions for the national policy contrast.

\begin{thebibliography}{9}
\bibitem{vivalt} E. Vivalt (2020), How Much Can We Generalize from Impact Evaluations?, \emph{Journal of the European Economic Association} 18(6), 3045--3089. Publisher abstract checked:
\url{https://doi.org/10.1093/jeea/jvaa019}.
\end{thebibliography}
\end{document}
